\section{Scientific rationality: the logical limits of a test}
\label{sec:p11-rat141-annex}

This proof makes explicit the elementary prediction/refutation distinction
in the selected chapter \parencite[pp.~747--748]{AndersenAndersen2021Scientific}.
It does not supply an inductive confirmation measure or a policy for
selecting which empirical premise to revise.

\begin{annexitem}{annex:p11.rat141.test-logic}{thm:p11-rat141-test-logic}{Boolean consequences and both countervaluations}
\begin{proof}
Take any valuation making \((H\land A)\to E\) and \(\neg E\)
true. It makes \(E\) false. If it made \(H\land A\) true,
the conditional would be false, contrary to the premise. Therefore
it makes \(H\land A\) false and \(\neg(H\land A)\) true.
This proves the first entailment by quantifying over every valuation
satisfying the premises.

If the additional premise \(A\) is true, the falsity of the conjunction
forces \(H\) false: otherwise both conjuncts would be true.
Thus \(\neg H\) is true in every valuation of the three premises,
which proves the second entailment.

The additional premise cannot generally be omitted. Set \(H\) true,
\(A\) false, and \(E\) false. The conjunction \(H\land A\)
is false, so the conditional is true; \(\neg E\) is also true,
but \(\neg H\) is false. This is a countervaluation to rejecting
the hypothesis from only a failed consequence and the conditional.

For the purported verification inference, set \(H\) false, \(A\)
true, and \(E\) true. The conditional is true because its antecedent
is false, and both observation and auxiliary are true, but \(H\)
is false. Hence those three premises do not entail the hypothesis.
The atoms' independent assignability makes each displayed valuation
available; the claim is invalidity of the general form, not of every
instance with further substantive premises.
\end{proof}
\end{annexitem}
